This study develops a $C^1$ Bogner--Fox--Schmit finite element formulation
for the coupled piezoelectric--flexoelectric response of a strain-gradient
strip with fractional mechanical stiffness memory. The coefficients of an
adopted positive local gradient energy are parametrized by the second-moment
tensor of a prescribed rectangular window. A Caputo Kelvin--Voigt law is
integrated by the L1 method in a reduced modal representation. The
implementation is assessed through closed-form electromechanical couplings,
mesh-refinement studies and limiting-case calculations. For the assumed
$\mathrm{BaTiO_3}$ parameters, the extensional coupling factor is
$k^2=0.5727$, compared with the closed-form value $0.5752$. Rotation of an
anisotropic window changes the computed coupling by approximately $14$--$17\,\%$.
When the strip thickness is reduced from \SI{10}{\micro\meter} to
\SI{1}{\micro\meter}, the flexoelectric bending coupling increases by
approximately a factor of $24$ in the archived mesh study, whereas the extensional piezoelectric
coupling changes by less than $0.08\,\%$. The first-period equivalent
decrement is $0.065$ at $\alpha=1$, reaches $0.105$ at $\alpha=0.6$,
and is $0.085$ at $\alpha=0.4$;
the damping also depends on which condensed stiffness operator carries the
memory. The computed design map spans flexoelectric coupling factors from
$3.6\times10^{-6}$ to $1.0\times10^{-2}$. Since the prescribed window
extends beyond the strip in part of the parameter range, these results
represent sensitivity of the adopted local gradient model rather than
material-calibrated device predictions. The calculations quantify the separate effects of window geometry,
electrical boundary conditions and fractional mechanical memory within
the prescribed model.
